\begin{center}
\SourceLeaf{045}\pdfbookmark[2]{29.}{article-029}29.
\end{center}
Quum quantitates auxiliares in ellipsi adhibitae $\varphi$, $E$, $M$ in hyperbola valores imaginarios obtineant, haud abs re erit, horum nexum cum quantitatibus realibus, quibus hic usi sumus, investigare: apponimus itaque relationes praecipuas, ubi quantitatem imaginarium $\BareRoot{-1}$ per $i$ denotamus.
\[
\sin\varphi=e=\frac{1}{\cos\psi}
\]
\[
\operatorname{tang}(45^0-\frac{1}{2}\varphi)=\BareRoot{\frac{1-e}{1+e}}=i\BareRoot{\frac{e-1}{e+1}}=i\operatorname{tang}\frac{1}{2}\psi
\]
\[
\operatorname{tang}\varphi=\frac{1}{2}\operatorname{cotang}(45^0-\frac{1}{2}\varphi)-\frac{1}{2}\operatorname{tang}(45^0-\frac{1}{2}\varphi)=-\frac{i}{\sin\psi}
\]
\[
\cos\varphi=i\operatorname{tang}\psi
\]
\[
\varphi=90^0+i\log(\sin\varphi+i\cos\varphi)=90^0-i\log\operatorname{tang}(45^0+\frac{1}{2}\psi)
\]
\[
\operatorname{tang}\frac{1}{2}E=i\operatorname{tang}\frac{1}{2}F=\frac{i(u-1)}{u+1}
\]
\[
\frac{1}{\sin E}=\frac{1}{2}\operatorname{cotang}\frac{1}{2}E+\frac{1}{2}\operatorname{tang}\frac{1}{2}E=-i\operatorname{cotang}F
\]
sive
\[
\sin E=i\operatorname{tang}F=\frac{i(uu-1)}{2u}
\]
\[
\operatorname{cotang}E=\frac{1}{2}\operatorname{cotang}\frac{1}{2}E-\frac{1}{2}\operatorname{tang}\frac{1}{2}E=-\frac{i}{\sin F}
\]
sive
\[
\operatorname{tang}E=i\sin F=\frac{i(uu-1)}{uu+1}
\]
\[
\cos E=\frac{1}{\cos F}=\frac{uu+1}{2u}
\]
\[
iE=\log(\cos E+i\sin E)=\log\frac{1}{u}
\]
sive
\[
E=i\log u=i\log\operatorname{tang}(45^0+\frac{1}{2}F)
\]
\[
M=E-e\sin E=i\log u-\frac{ie(uu-1)}{2u}=-\frac{iN}{\lambda}\raisebox{.45ex}{.}
\]
Logarithmi in his formulis sunt hyperbolici.
